\section{Embedding $\mathbb{W}$ into $\mathbb{Z}$}
\label{sec:embedding-w-into-z}

\begin{topicbox}{Embedding $\mathbb{W}$ into $\mathbb{Z}$}
\begin{exposition}
This section fixes the canonical embedding of $\mathbb{W}$ into $\mathbb{Z}$ and records that it is injective, structure-preserving, and order-reflecting. The map is the bridge that lets elements of $\mathbb{W}$ be regarded as elements of $\mathbb{Z}$ after identification.
\end{exposition}
\end{topicbox}
\begin{tcolorbox}[colback=defbox, colframe=defborder, arc=2pt,
  left=6pt, right=6pt, top=4pt, bottom=4pt,
  title={\small\textbf{Definition (Embedding $\mathbb{W}$ into $\mathbb{Z}$)}},
  fonttitle=\small\bfseries]
\begin{definition}[Embedding $\mathbb{W}$ into $\mathbb{Z}$]
\label{def:embedding-w-into-z}
The \emph{canonical embedding} $\iota:\mathbb{W}\to \mathbb{Z}$ is defined by
\[
\iota(w)=[(w,0)] \quad(\text{the difference class of } w \text{ and } 0).
\]
\end{definition}
\end{tcolorbox}
\begin{remark*}[Standard quantified statement]
\[
\iota:\mathbb{W}\to \mathbb{Z},\qquad \iota(w)=[(w,0)] \quad(\text{the difference class of } w \text{ and } 0).
\]
\end{remark*}
\begin{remark*}[Interpretation]
The map sends each element of $\mathbb{W}$ to its natural representative in $\mathbb{Z}$. The next result certifies that this representative choice preserves all arithmetic and order, so $\mathbb{W}$ may be identified with its image in $\mathbb{Z}$.
\end{remark*}
\begin{dependencies}
\begin{itemize}
  \item TODO: link to the construction of $\mathbb{Z}$ as difference classes of whole numbers.
  \item \hyperref[def:number-system-embedding]{Embedding of Number Systems}
\end{itemize}
\end{dependencies}
\begin{tcolorbox}[colback=thmbox, colframe=thmborder, arc=2pt,
  left=6pt, right=6pt, top=4pt, bottom=4pt,
  title={\small\textbf{Theorem ($\mathbb{W}$ Embeds in $\mathbb{Z}$)}},
  fonttitle=\small\bfseries]
\begin{theorem}[$\mathbb{W}$ Embeds in $\mathbb{Z}$]
\label{thm:w-embeds-in-z}
The canonical map $\iota:\mathbb{W}\to \mathbb{Z}$ is an embedding of ordered arithmetic systems: it is injective, preserves $0$, $1$, addition, and multiplication, and reflects the order.
\hyperref[prf:w-embeds-in-z]{Go to proof.}
\end{theorem}
\end{tcolorbox}
\begin{remark*}[Standard quantified statement]
\[
\iota:\mathbb{W}\to \mathbb{Z} \text{ is injective, structure-preserving, and order-reflecting.}
\]
\end{remark*}
\begin{remark*}[Interpretation]
Once this theorem is proved, the identification theorem of \S\ref{sec:embeddings-as-structure-preserving-maps} applies, and $\mathbb{W}$ may be treated as the substructure $\iota(\mathbb{W})\subseteq \mathbb{Z}$.
\end{remark*}
\begin{dependencies}
\begin{itemize}
  \item TODO: link to the construction of $\mathbb{Z}$ as difference classes of whole numbers.
  \item \hyperref[def:embedding-w-into-z]{Embedding $\mathbb{W}$ into $\mathbb{Z}$}
  \item \hyperref[thm:embedding-identifies-with-image]{Embeddings Identify a System with Its Image}
\end{itemize}
\end{dependencies}
